\documentclass[11pt,a4paper]{article}

% ========================================
% PACKAGES
% ========================================
\usepackage[utf8]{inputenc}
\usepackage[T1]{fontenc}
\usepackage{lmodern}
\usepackage[english]{babel}
\usepackage[margin=2.5cm]{geometry}

% Mathematics
\usepackage{amsmath, amssymb, amsthm}
\usepackage{mathtools}

% Graphics and colors
\usepackage{graphicx}
\usepackage{xcolor}
\usepackage{tikz}
\usetikzlibrary{positioning, arrows.meta, shapes.geometric, calc}
\usepackage[most]{tcolorbox}

% Lists and formatting
\usepackage{enumitem}
\usepackage{parskip}
\usepackage{fancyhdr}
\usepackage{multirow}

% Code listings
\usepackage{listings}

% Hyperlinks
\usepackage{hyperref}
\hypersetup{
    colorlinks=true,
    linkcolor=blue,
    urlcolor=blue,
    citecolor=blue
}

% ========================================
% THEOREM ENVIRONMENTS
% ========================================
\theoremstyle{definition}
\newtheorem{definition}{Definition}[section]
\newtheorem{example}{Example}[section]
\newtheorem{exercise}{Exercise}[section]

\theoremstyle{plain}
\newtheorem{theorem}{Theorem}[section]
\newtheorem{lemma}[theorem]{Lemma}
\newtheorem{proposition}[theorem]{Proposition}
\newtheorem{corollary}[theorem]{Corollary}

\theoremstyle{remark}
\newtheorem*{remark}{Remark}
\newtheorem*{observation}{Observation}

% ========================================
% TABLE OF CONTENTS DEPTH
% ========================================
\setcounter{tocdepth}{2}

% ========================================
% HEADER AND FOOTER
% ========================================
\setlength{\headheight}{14pt}
\pagestyle{fancy}
\fancyhf{}
\lhead{\leftmark}
\rhead{CC}
\cfoot{\thepage}
\renewcommand{\sectionmark}[1]{\markboth{#1}{}}

% ========================================
% DOCUMENT INFORMATION
% ========================================
\title{\textbf{Cybercrime}\\
\large Artificial Intelligence}
\author{Jacopo Parretti}
\date{II Semester 2025-2026}

% ========================================
% DOCUMENT
% ========================================
\begin{document}

\maketitle
\tableofcontents
\newpage


% ========================================
\part{Crime in the Digital Society}
% ========================================

\section{Digitalisation: Definition and Short History}

\subsection{Digital Data}

\begin{definition}[Digital data]
Electronic representation of information in a format or language that computers can read and understand. In more technical terms: information in binary format (0--1) that is converted into a digital format readable by a computer.
\end{definition}

\subsection{A Brief History of Digitalisation}

\begin{itemize}
    \item \textbf{1663}: John Graunt applies statistical analysis to mortality data from the bubonic plague --- an early precursor to data-driven reasoning.
    \item \textbf{1865}: R. Millar Devens coins the concept of business intelligence.
    \item \textbf{1884}: Herman Hollerith builds the tabulating machine for the US Census; later founds what becomes IBM.
    \item \textbf{1943}: Colossus, the first programmable electronic computer, is used to decipher Nazi communications at Bletchley Park.
    \item \textbf{1989}: Tim Berners-Lee creates the World Wide Web.
    \item \textbf{1995}: First supercomputer.
\end{itemize}

\subsection{From Digital Data to Big Data}

\begin{itemize}
    \item \textbf{Since 1997}: Emergence of social networks.
    \item \textbf{1999}: The term \textbf{Internet of Things (IoT)} enters use --- the extension of the Internet to the world of physical objects.
    \item \textbf{2005}: Roger Mougalas coins the term \textbf{Big Data}: a set of digital data so large that it cannot be analysed and stored using traditional tools; specific technologies and analytical methods are needed to extract value from it.
\end{itemize}


\section{Case Study: Smart Toys and Criminal Risks}

This case study illustrates how IoT consumer devices --- specifically smart toys --- introduce concrete criminal risks by collecting and transmitting sensitive data about children.

\subsection{The Cayla Doll}

The Cayla doll was an internet-connected toy capable of recording and transmitting children's voice data. In 2017, Germany banned the device over fears that hackers could exploit it to target children. The case raises a fundamental question that recurs throughout the course: \textit{unfounded fears or documented reality?}

\begin{tcolorbox}[colback=blue!5!white,colframe=blue!60!black,title={Reference}]
Justin Huggler, ``Germany bans internet-connected dolls over fears hackers could target children'', \textit{The Telegraph}, 17 February 2017.
\end{tcolorbox}

\subsection{CloudPets}

CloudPets were internet-connected stuffed animals that allowed children and parents to exchange voice messages via a companion app. A security analysis by Troy Hunt revealed that the backend database --- containing millions of voice recordings and user credentials --- was left exposed on the internet without authentication, was subsequently accessed by unauthorised parties, and the data was held for ransom. This is a documented case of a data breach compounded by ransom, affecting a particularly vulnerable population.

\subsection{Wiggy}

Wiggy was a crowdfunded ``smart piggy bank'' that linked to a child's savings account via a mobile app, allowing parents and relatives to deposit money remotely. The device combined financial account access with an IoT object marketed at children, creating a direct bridge between physical play and banking infrastructure --- a surface area for financial fraud and account compromise.


\section{Crime in the Digital Society}

The digital environment does not create crime \textit{ex novo}, but it transforms its scale, speed, authorship, and victim profile. The following dimensions illustrate this transformation through concrete examples.

\subsection{Crime Types}

Digital infrastructure enables new crime typologies or drastically scales existing ones:

\begin{itemize}
    \item \textbf{Cyberattacks on critical infrastructure}: In February 2023, a pro-Russian hacker group targeted Italian government websites and corporate portals (Carabinieri, Bper bank, A2A, TIM) coinciding with Prime Minister Meloni's visit to Kyiv. The Italian National Cybersecurity Agency (ACN) managed to mitigate the effects.
    \item \textbf{Election manipulation}: Investigative reporting identified a network led by ``Tal Hanan'' (former Israeli special forces operative) that allegedly manipulated 30 elections worldwide by deploying hacking, disinformation, and social engineering campaigns on behalf of paying clients.
    \item \textbf{Online fraud and sextortion}: Italian police data and news reports document a continuous rise in online fraud, with sextortion and fake social profiles representing a significant and growing share.
    \item \textbf{AI-assisted fraud}: In February 2025, a fraud scheme used AI-generated voice cloning to impersonate Italian Defence Minister Crosetto and a cloned ministerial phone number to solicit money from major Italian entrepreneurs (Armani, Moratti, Tronchetti Provera, Della Valle) under the pretext of ``freeing hostages''. At least one victim paid.
\end{itemize}

\subsection{Crime Authors}

The digital environment radically diversifies who can commit crime and at what scale:

\begin{itemize}
    \item \textbf{Individual actors}: Cody Wilson (b.\@ 1988), founder of Defense Distributed, published open-source gun designs (``wiki weapons'') for 3D printing online, gaining international notoriety in 2013 with the Liberator --- the first widely available 3D-printed pistol.
    \item \textbf{Teenage actors}: Graham Ivan Clark, at age 17, hacked Twitter's internal admin tool in 2020 and took control of high-profile accounts (Obama, Gates, Musk, West) to run a Bitcoin scam that defrauded victims of over \$100,000. He pleaded guilty to a three-year sentence.
    \item \textbf{Nation-state actors}: Bureau 121, North Korea's internal hacking unit, has been linked to several major malware events, most notably the \textbf{WannaCry ransomware attack} (2017), which infected around 300,000 devices across 150 countries and caused an estimated \$4 billion in damages. The US sanctioned North Korea over WannaCry in 2019.
    \item \textbf{AI systems as unintended actors}: Microsoft's Tay chatbot (2016) was corrupted by coordinated users within 24 hours of launch and began producing racist and misanthropic content at scale --- a case where the crime author is ambiguous and the platform's own technology becomes the vector.
\end{itemize}

\subsection{Crime Victims}

Victimisation in the digital society is broader, faster, and less visible than in traditional crime:

\begin{itemize}
    \item \textbf{Sextortion}: FBI data documents a troubling increase in sextortion cases targeting minors. In one case (May 2022), a 17-year-old died by suicide hours after being scammed.
    \item \textbf{Mass data breach}: The Ashley Madison hack exposed the personal data of 32 million users of an extramarital affairs platform, with ongoing blackmail and reputational damage for victims years after the breach.
    \item \textbf{Cyber kidnapping}: A new criminal trend targeting foreign exchange students and their families. Police found a Chinese exchange student isolated and frightened in a tent in the Utah mountains; the perpetrators had convinced the victim to stage a kidnapping while extorting the family remotely.
    \item \textbf{Digital-enabled physical crime}: Thieves monitored Facebook posts to identify when homes were empty before burglarising them. In the Kim Kardashian Paris hotel robbery (2016), the perpetrators reportedly used her social media to track her location and movements.
    \item \textbf{Digital humiliation}: In Milan (2023), a perpetrator filmed and shared on social media footage of a physical assault she committed, leading victims' counsel to demand a public apology on TikTok as part of the settlement.
\end{itemize}

\subsection{Social Control and Policing}

Digital technologies also reshape the instruments of social control:

\begin{itemize}
    \item \textbf{Facial recognition}: In 2020, the European Commission considered a five-year ban on facial recognition in public spaces, highlighting the tension between law enforcement utility and civil liberties.
    \item \textbf{Algorithmic risk assessment}: Police forces increasingly use algorithmic tools to assess threat levels and predict criminal behaviour. A 2016 ProPublica investigation (``Machine Bias'') documented that the COMPAS recidivism prediction software used across US courts exhibited systematic racial bias, falsely flagging Black defendants as future criminals at nearly twice the rate of white defendants.
\end{itemize}

\begin{tcolorbox}[colback=yellow!5!white,colframe=orange!80!black]
A recurring tension throughout the course: the same digital tools that enable crime also constitute the primary instruments of crime detection and social control. This duality --- and its legal, ethical, and criminological implications --- is the central analytical lens of the course.
\end{tcolorbox}


% ========================================
\part{Examples of Criminological Studies}
% ========================================

\section{Confessions of a Migrant Smuggler}

An ethnographic study by Andrea Di Nicola and Giampaolo Musumeci into migrant smuggling as a classic organised crime activity with first elements of digitalisation. The research was conducted in the years before the Mediterranean migration crisis and consisted of a two-year journey inside the minds of human smugglers, presenting the downside of the coin and debunking stereotypical ideas about ``scafisti''.

\subsection{First Elements of Digitalisation}

Smugglers were found to use technology and digital tools for several purposes:

\begin{itemize}
    \item For organisational purposes and to conduct their activities.
    \item To communicate with other smugglers in the network.
    \item To recruit clients and market their services.
    \item To acquire the tools necessary for the business (e.g.\ falsified documents, credit card information).
\end{itemize}

New means of communication adopted by smuggling networks include Skype, satellite phones and cellphones, wireless networks, email, and instant messaging apps.


\section{Project Surf and Sound: The Digitalisation of Migrant Smuggling}

The Surf and Sound project aimed at ``improving and sharing knowledge on the Internet's role in the human trafficking process''. It was the first research project to produce systematic knowledge on the use of the Internet (including social networks) in trafficking in human beings (THB) and migrant smuggling, disseminating findings through workshops for law enforcement agencies (LEAs) and NGOs. The project was co-financed by the European Commission --- DG Home Affairs within the ISEC 2013 ``Prevention and Fight against Crime -- Trafficking in Human Beings'' programme.

\subsection{Methodology}

The study employed a mixed-methods design:

\begin{itemize}
    \item \textbf{Internet exploration (virtual ethnography)}: passive and active virtual ethnography ($N = 939$).
    \item \textbf{In-depth interviews} ($N = 58$): LEAs (16), NGOs (13), victims/clients (14), traffickers/smugglers (13), and cybercrime experts (2).
    \item \textbf{Online web survey} ($N = 995$): analysing the experience of internet users with web content linked to human trafficking and people smuggling.
\end{itemize}

\subsection{Key Results}

\begin{enumerate}
    \item The Internet plays a relevant role in the smuggling of migrants.
    \item Content linked to migrant smuggling services on social media websites is easily accessible.
    \item Social media websites facilitate the advertisement of smuggling services; Facebook and Instagram are more used than Twitter.
    \item Social media pages and profiles run by smugglers provide information on travels.
    \item Social media profiles and pages run by smugglers are very effective and responsive.
    \item Smugglers provide potential migrants with detailed information via mobile applications.
    \item Social media websites play a key role in the decision-making process of potential migrants.
    \item Social media websites are used by migrants to post feedback about smuggling services.
    \item The Internet is used also during and after the journey.
    \item The deep/dark web seems to play a marginal role.
\end{enumerate}


\section{Project FAKECARE: The Digitalisation of the Illegal Trade of Medicines}

\subsection{The Problem}

An estimated \textbf{10\%} of medicines worldwide are fake, illegal, or substandard. There are significant geographical differences: in developed countries the figure is below 1\%, while in developing countries it ranges from 15--30\% (sometimes reaching 50\%). This results in an estimated 100,000 to 700,000 deaths annually, with a global estimated turnover of 75 billion USD per year (\textbf{10 billion EUR/year in the EU}).

\subsection{Falsified Medicinal Products}

The use of several terms (``counterfeit'', ``falsified'', ``false'', ``tampered'') creates great confusion. For clarity, falsified medicinal products are \textbf{medicines that are dangerous} for private and public health. The umbrella category of \textbf{illegally traded (or produced) medicines} (ITMs) includes:

\begin{itemize}
    \item Falsified branded medicines
    \item Falsified generic medicines
    \item Illegal medicines
    \item Medicines without authorisation
    \item Medicines with expired authorisation
    \item Legal (and genuine) medicines sold without a valid prescription
\end{itemize}

\subsection{Methodology}

\paragraph{Quantitative methods.}
\begin{itemize}
    \item \textbf{Web survey}: profiling the risk for customers to buy a falsified medicinal product online.
    \item \textbf{Honey-pot websites}: two fake online pharmacy websites set up to attract customers, automatically monitor their behaviour, and calculate statistics on visitors, clicks, and checkouts.
    \item \textbf{Quantitative content analysis}: on online pharmacies, customers' reviews on illegal online pharmacies, and news from the press.
    \item \textbf{Statistics} on the online trade of falsified medicines using data provided by regulatory agencies, law enforcement agencies, and Interpol.
\end{itemize}

\paragraph{Qualitative methods.}
\begin{itemize}
    \item Analysis of \textbf{EU regulations} in the field of the online trade of falsified medicines.
    \item \textbf{Virtual ethnography}: exploring online forums and communities (including social networks).
    \item \textbf{Crime script analysis}: analysing investigative and judicial cases, news from the press.
    \item \textbf{Critical quantitative analysis of secondary sources}.
\end{itemize}

\subsection{Crime Script of the Online Trade}

The online trade can be seen as a process rather than a single event, involving stages in which resources and locations are required and decisions are made. The focus is mainly on the \textbf{supply side}:

\begin{enumerate}
    \item \textbf{Preparation}: identification of the criminal opportunities to commit the online trade in falsified medicinal products.
    \item \textbf{Pre-activity}: logistical or transactional instrumental actions prior to the activity.
    \item \textbf{Activity}: concrete actions necessary to carry out the trade, which can take the form of various offences of non-compliance.
    \item \textbf{Post-activity}: logistical or transactional activities necessary to leave the scene.
\end{enumerate}

\subsection{The Role of the Internet in Distribution}

There are significant geographical differences in distribution channels:

\begin{itemize}
    \item \textbf{Developing countries}: legal distribution chain.
    \item \textbf{Developed countries}: illegal distribution chain (illegal websites, online marketplaces, social media websites, etc.).
\end{itemize}

An estimated \textbf{94\%} of online pharmacies (around 34,000) operate illegally.

\subsection{Indicators of an Illegal Online Pharmacy}

\begin{itemize}
    \item No request of a prescription for prescription-only medicines
    \item Low prices, promotions, and offer of trial packs
    \item Use of different languages, grammar mistakes
    \item Absence of a physical address
    \item Promotion of medicinal products on the home page
    \item Great attention to providing anonymity
    \item Use of ``testimonials''
\end{itemize}

\subsection{Distribution Channels}

Falsified medicines are distributed through multiple channels: social media platforms (Facebook, Twitter, Instagram), online marketplaces (e.g.\ Alibaba), and the dark web (e.g.\ the former Silk Road marketplace).

\subsection{EU Common Logo for Online Sale of Medicines}

As of 1 July 2015, all legally operating online pharmacies/retailers in the EU must display a common logo. The logo must be clearly displayed on every page of the website and links to the website of the national competent authority and to the national list of legally operating online pharmacies/retailers.

However, the system is \textbf{vulnerable to fraudsters}. A vulnerability assessment demonstrated that the EU common logo and the verification scheme can be counterfeited in a few minutes, with basic computer skills, spending only the money necessary to purchase a fraudulent domain. The steps are:

\begin{enumerate}
    \item Copy the logo from a legitimate website.
    \item Copy the institutional webpage of the national authority.
    \item Falsify the institutional webpage to list the illegal pharmacy.
    \item Purchase a fraudulent domain with a similar URL.
\end{enumerate}

In addition, the verification scheme has not been well implemented in some cases: the logo is not displayed at all, or the link does not work or points to pages other than the national register.


% ========================================
\part{Cybercrime Typologies: Cyber-Dependent Crimes}
% ========================================

\section{Categories of Cybercrime}

The term \textbf{cybercrime} covers a wide range of deviant and criminal behaviours that are constantly changing and evolving. It is unlikely that there is an exhaustive list or universally accepted definition of cybercrime. More common definitions refer to categorisations of cybercrime (Phillips et al., 2022), and these categories are not rigid --- there is the possibility of partial overlap. Cybercrime is often defined using \textbf{dichotomies and trichotomies}.

\subsection{The Brenner--Wall Trichotomy}

Brenner (2007) proposed a dichotomy later extended by Wall (2010) into a trichotomy:

\begin{itemize}
    \item \textbf{Cyber-dependent crime}: crimes that arose with technology and cannot exist outside the digital world.
    \item \textbf{Cyber-enabled crime}: the process of digitalisation has facilitated the commission of these crimes.
    \item \textbf{Cyber-assisted crime}: traditional crimes in which the use of the computer is not essential but rather incidental.
\end{itemize}

\subsection{The Gordon--Ford and Sarre--Lau--Chang Trichotomy}

Gordon and Ford (2006) distinguished:

\begin{itemize}
    \item \textbf{Type I}: technical cybercrimes.
    \item \textbf{Type II}: cybercrimes that involve human contact.
\end{itemize}

Sarre, Lau, and Chang (2018) added a third category:

\begin{itemize}
    \item \textbf{Type III}: crimes perpetuated by AI, robots, bots, and self-learning technology.
\end{itemize}


\section{Cyber-Dependent Crimes}

\begin{definition}[Cyber-dependent crime]
High-tech crimes that can only be committed using computers, computer networks, or other forms of ICT (Wall, 2005; Europol, 2017). These are crimes that \textbf{can only be committed in cyberspace} and would not exist without the Internet or technology. They represent a new form of crime that has emerged with the digital society --- also referred to as \textit{true cyber crime} or \textit{real cyber crime}.
\end{definition}

\subsection{Subcategories by Mode of Execution}

Based on how they are carried out, cyber-dependent crimes can be divided into \textbf{two subcategories} (McGuire \& Dowling, 2013):

\begin{enumerate}
    \item \textbf{Illicit intrusions into computer networks}:
    \begin{itemize}
        \item Hacking
        \item Unauthorised access to computer networks, computers, and mobile devices exploiting security gaps and vulnerabilities in order to collect personal data and useful information, or to deface websites and launch DoS/DDoS attacks
    \end{itemize}
    \item \textbf{Interruption or damage to the functioning of computers and networks}:
    \begin{itemize}
        \item Malware
        \item Spam
        \item DoS
        \item DDoS
        \item Botnets
    \end{itemize}
\end{enumerate}

\subsection{Subcategories by CIA Triad}

Brar and Kumar (2018) proposed a \textbf{threefold definition} based on the pillars of cybersecurity:

\begin{itemize}
    \item \textbf{Attacks on confidentiality}: eavesdropping, traffic analysis, password attacks, snooping, keyloggers, social engineering.
    \item \textbf{Attacks on integrity}: salami attacks, data diddling, cross-site scripting (XSS), SQL injection, session hijacking.
    \item \textbf{Attacks on availability}: DoS, DDoS.
\end{itemize}


\section{Case Studies}

\subsection{Ransomware Attacks on Hospitals}

On 3 August 2023, a ransomware attack hit a California-based healthcare system operating 16 hospitals and more than 165 clinics and centres across four US states. As a result, several services were shut down, facilities closed, patients and ambulances diverted, and digital records were unavailable. Cyberattacks on hospitals are on the rise worldwide (Argaw, Bempong, Eshaya-Chauvin et al., 2019).

\subsection{Ransomware Attacks on Law Enforcement}

In August 2023, a data breach at a company that prints ID cards and lanyards compromised data from customers including Greater Manchester Police and the London Metropolitan Police. The exposed data included names, ranks, photos, and serial numbers. Europol conducted an operation that led to the arrest of a member of the Ragnar Locker group and the seizure of its infrastructure.

\subsection{DDoS Attacks}

DDoS attacks are one of the biggest cyber threats and have evolved to avoid detection and increase damage (Bergamini de Neira, Kantarci, Nogueira, 2023). In October 2023, a record-breaking DDoS attack was reported against Google, Cloudflare, and AWS, peaking at over 398 million requests per second (RPS). Amazon (June 2020) and Google (2022) had previously reported stopping the then-largest DDoS attacks. The scale has grown from 26 to 398 million requests per second.

\subsection{Ashley Madison}

In July 2015, a hacker collective called ``The Impact Team'' infiltrated Ashley Madison --- a dating website for those seeking extramarital affairs --- stealing over 60 gigabytes of sensitive data (user profiles, financial transactions, sexual orientations, internal communications). Personal information of approximately 32 million users was exposed, leading to personal crises, suicides, and allegations of fraud and deceptive business practices by the platform.

\subsection{The Role of Artificial Intelligence}

AI introduces new dimensions to cyber-dependent crime:

\begin{itemize}
    \item \textbf{Acoustic side-channel attacks}: researchers developed deep-learning systems capable of hacking passwords using key acoustics with over 90\% accuracy (Harrison, Toreini, Mehrnezhad).
    \item \textbf{AI-enhanced social engineering}: chatbots such as ChatGPT can be exploited to create convincing phishing emails by correcting spelling and grammatical errors that would otherwise alert human readers. ``ChatGPT may therefore provide criminals with new opportunities, particularly for crimes involving social engineering'' (Europol, 2023).
\end{itemize}



% ========================================
\part{Cybercrime Typologies: Cyber-enabled Crimes}
% ========================================

\section{Introduction and Definition}

The digital society has given a \textbf{strong boost to a number of traditional crimes} that have been enhanced by the use of computers, computer networks, or other forms of information and communication technology (ICT). In contrast to cyber-dependent crime, \textbf{cyber-enabled crime} (McGuire \& Dowling, 2013) \textbf{can also be committed without the use of the digital world}: it is not born with technology, but technology transforms its scale, reach, and ease of commission.

A key distinction: cyber-enabled crime \textit{has always existed}, but it is now more \textbf{ubiquitous}, \textbf{faster}, and more \textbf{global}. The internet does not create new criminal motivations, but it dramatically expands opportunities and reduces barriers.

\begin{definition}[Cyber-enabled crime]
Traditional crimes for which the Internet has opened up entirely new possibilities and opportunities; crimes that are facilitated or have been made easier (i.e., enabled) by cyber technology (Phillips et al., 2022). These are crimes where the computer plays a role but that could still be committed without the involvement of the computer. Cyber-enabled crimes range from white-collar crime to drug trafficking, to online harassment, terrorism, and beyond.

Also referred to as:
\begin{itemize}
    \item \textbf{Hybrid crimes} (Wall, 2005)
    \item \textbf{Computer-related crimes} (Broadhead et al., 2018)
    \item \textbf{People crimes} (Black et al., 2019)
\end{itemize}
\end{definition}

\begin{remark}
The boundary between cyber-dependent and cyber-enabled crime is not always clear-cut. The categories are analytically useful but descriptively fluid: the same conduct can shift categories depending on the specific modus operandi employed.
\end{remark}


\section{Subcategories of Cyber-enabled Crime}

Cyber-enabled crimes are organised into two broad clusters based on the nature of the harm caused:

\subsection{Cyber-enabled Crimes Related to Economy}

This cluster includes offences in which the primary motive is financial gain and the digital environment provides the infrastructure, anonymity, and scale needed to commit them:

\begin{itemize}
    \item Fraud (e-commerce fraud, mass-marketing frauds, phishing scams, pharming scams, ``online romance'' frauds)
    \item Intellectual property crimes
    \item Sale of illegal items on online marketplaces (including the deep or dark web)
    \item Cyber-laundering
\end{itemize}

\subsection{Cyber-enabled Crimes Related to Interpersonal Violence}

This cluster encompasses offences targeting individuals, often with psychological, reputational, or sexual harm as the primary outcome:

\begin{itemize}
    \item Cyber-bullying
    \item Cyber-mobbing
    \item Cyber-stalking
    \item Cyberharassment
    \item Sextortion
    \item Online grooming
    \item Trolling
    \item Griefing
    \item Assisted cyber suicide
    \item Image-based abuse (IBA)
    \item Production and dissemination of child pornography
\end{itemize}


\section{Cyber-enabled Crimes Related to Economy}

\subsection{The Industrialisation of Fraud}

The world of fraud is one of the areas that has benefited most from the introduction of the internet. The digital society has \textbf{industrialised fraud}: it lowered the barriers to entry, allowed offenders to reach exponentially more victims simultaneously, and reduced the risk of detection and punishment. Digital society is therefore particularly criminogenic in terms of fraud --- so much so that one could coin the adjective ``\textbf{fraudogenic}'' (Button \& Cross, 2017).

Characteristics that make online fraud structurally different from traditional fraud:
\begin{itemize}
    \item \textbf{Scale}: a single offender can target thousands of victims simultaneously.
    \item \textbf{Anonymity}: digital identities are easy to fabricate and difficult to trace.
    \item \textbf{Globalisation}: victims and perpetrators can be in entirely different jurisdictions, complicating prosecution.
    \item \textbf{Low detection risk}: compared to in-person fraud, the likelihood of identification and punishment is significantly reduced.
\end{itemize}

\subsection{Fraud Typologies}

McGuire and Dowling (2013) identify the following typologies of online fraud:

\begin{itemize}
    \item \textbf{Online frauds and e-commerce frauds}: broad category covering any deceptive scheme conducted via the internet or digital commerce platforms. Includes misrepresentation of goods, non-delivery of purchased items, and payment fraud.

    \item \textbf{Fraudulent sales}: carried out via online auction or retail sites, or through entirely bogus websites created specifically to defraud buyers. The seller collects payment but delivers counterfeit, inferior, or no goods at all.

    \item \textbf{Online mass marketing scams}: individuals are tricked into paying money upfront --- for example, to help someone, to participate in a business opportunity, or to claim a supposed prize --- with the promise that they will receive a larger sum of money in return at a later date. Classic examples include advance-fee fraud (``419 scams'') and lottery scams.

    \item \textbf{Phishing scams}: a user tricks the victim into providing personal information, financial details, or access codes by posing as a reputable entity or person in an email or other form of communication. The communication typically mimics the branding and language of legitimate institutions (banks, government agencies, e-commerce platforms).

    \item \textbf{Pharming scams}: a user is directed to a fake website --- sometimes via phishing emails, sometimes through DNS poisoning --- to enter their personal information. Unlike phishing, pharming can redirect users even when they type the correct URL into their browser.

    \item \textbf{Online romance scams}: the victim is contacted via social networks or dating websites and subjected to a prolonged online `relationship', at the end of which they are tricked into handing over personal details or large sums of money. Perpetrators typically fabricate elaborate identities (military personnel stationed abroad, wealthy professionals) and cultivate trust over weeks or months.
\end{itemize}

\subsection{Beyond Fraud: Other Economic Cyber-enabled Crimes}

\begin{itemize}
    \item \textbf{Intellectual property crimes}: product piracy, counterfeiting, forgery, and the sale of counterfeit goods enabled by global online marketplaces. Digital distribution has dramatically lowered the cost of distributing pirated software, films, music, and physical goods.

    \item \textbf{Sale of illegal items on online marketplaces}: the internet, and particularly the deep and dark web, hosts marketplaces where users can find and purchase weapons, drugs, human organs, stolen data, hacking tools, and a variety of illegal services. The dark web provides layers of anonymity through tools such as Tor that make both buyers and sellers difficult to identify.

    \item \textbf{Cyber laundering}: criminals can clean the proceeds of illegal activities using cyber laundering systems. Some or all stages of the classical money laundering process (placement, layering, integration) may take place in cyberspace --- through cryptocurrency transactions, online gambling platforms, or fictitious e-commerce sales --- in order to erase the criminal origin of assets.
\end{itemize}


\section{Cyber-enabled Crimes Related to Interpersonal Violence}

\subsection{Malicious and Offensive Communications}

\begin{itemize}
    \item \textbf{Cyber-bullying}: intimidation and bullying of young people (less often children) through ICTs, including the creation of special websites with material about the person being persecuted, or by posting offensive reviews and comments. Cyber-bullying can be carried out by individuals or by coordinated groups and often extends into the physical world, affecting the victim's school and social life.

    \item \textbf{Cyber-mobbing}: psychological violence in the form of bullying that uses electronic communication to insult, threaten, or humiliate a victim. Unlike cyber-bullying (which tends to target peers in school settings), cyber-mobbing is typically carried out by a group of people who collectively bully one of their group members or any user, often in workplace or community contexts.
\end{itemize}

\subsection{Online Violence Against Other Individuals}

\begin{itemize}
    \item \textbf{Cyber-stalking}: the systematic pursuit of a person, a group of people, or an organisation via the internet with the aim of intimidating them, blackmailing them, and making demands. It involves repeated, targeted conduct that causes fear or distress in the victim. Digital tools allow stalkers to monitor victims' movements, contacts, and routines at a level of detail previously impossible.

    \item \textbf{Cyber-harassment}: a repeated psychological cyberattack through persistent insults, harassment, demands, and virtual aggression. It typically manifests as a continuous flow of numerous messages, intrusive calls around the clock, and conversations of a humiliating and insulting nature. The persistence and visibility of online harassment can amplify harm well beyond what equivalent in-person conduct would cause.

    \item \textbf{Image-based abuse (IBA)}: a continuum of different forms of sexual violence that includes:
    \begin{itemize}
        \item Non-consensual intimate image (NCII) distribution
        \item Revenge porn
        \item Deepfakes or fake pornography (AI-generated sexual imagery using real individuals' likenesses)
        \item Sextortion (the use or threatened use of intimate images to coerce victims into further sexual acts or payments)
    \end{itemize}
\end{itemize}

\subsection{Online Sexual Offences Against Children}

\begin{itemize}
    \item \textbf{Online grooming}: a form of sexual abuse of children and commercial sexual exploitation in which adult users employ digital platforms to establish trust with underage victims, with the ultimate intention of luring them into sexual contact or exploitation. Grooming typically unfolds over time, with the offender progressively desensitising the victim to sexual content.

    \item \textbf{Production and distribution of child pornography}: online platforms, including the dark web, are used to produce, share, and trade child sexual abuse material (CSAM). The digital environment has enabled the formation of large, international networks of offenders who share material and coordinate offending.
\end{itemize}

\subsection{Other Forms of Cyber-enabled Interpersonal Crime}

\begin{itemize}
    \item \textbf{Trolling}: social provocation in which virtual conversation partners are deliberately drawn into controversy in order to stir up a scandal or inflict emotional distress in network communication. Trolling ranges from low-level harassment to coordinated campaigns targeting individuals.

    \item \textbf{Griefing}: the deliberate infliction of moral and material damage on participants in online computer games --- sabotaging gameplay, destroying items, or psychologically tormenting other players. While often considered a minor infraction, in some cases it shades into extortion or fraud.

    \item \textbf{Assisted cyber suicide}: a form of psychological abuse in which the victim is driven to suicide by psychological manipulation, sustained pressure, or encouragement on the internet. It sits at the intersection of cyber-harassment and criminal facilitation of suicide.
\end{itemize}


\section{Case Studies}

\subsection{E-commerce Fraud: Belgian VAT Scheme (March 2023)}

In March 2023, Operation ``Silk Road'' uncovered a large-scale VAT fraud scheme involving Chinese e-commerce exporters. The perpetrators evaded Belgian import taxes by falsely declaring that consumer goods (electronics, toys, accessories) were destined for other EU Member States, thereby qualifying for a VAT exemption. The scheme relied on a network of bogus companies registered in France, Germany, Hungary, Italy, Poland, and Spain, supported by false invoices, forged transport documents, and fictitious identities. The total amount of VAT fraud was estimated at \textbf{303 million euros}. The case illustrates how cross-border digital trade creates systematic vulnerabilities in tax enforcement regimes, and how international criminal networks can exploit regulatory gaps between Member States.

\subsection{Cyber Laundering: Estonian Cryptocurrency Fraud (November 2022)}

In November 2022, two Estonian nationals were arrested in Tallinn for their involvement in a multi-layered fraud and money laundering scheme centred on cryptocurrencies. The perpetrators operated on two fronts: they tricked victims into signing fraudulent equipment rental contracts, and separately induced victims to invest in a fictitious virtual currency bank. In total, victims paid more than \textbf{\$575 million} to the defendants' companies. The proceeds were then laundered through shell companies and used to purchase real estate and luxury cars in Estonia. The case demonstrates the intersection of classic investment fraud with cyber laundering, and the use of cryptocurrency's opacity to obscure illicit financial flows.

\subsection{Online Romance Fraud (March 2023, Italy)}

In March 2023, Italian police arrested eight individuals for serious fraud, money laundering, and counterfeiting, following an investigation into a romance scam network. Victims (primarily women) were contacted on Facebook by a fake profile named ``Larry Brooks'', a self-proclaimed US Army officer stationed in Syria. Over the period 2018--2021, the criminal group defrauded \textbf{32 victims} of approximately \textbf{\texteuro 400,000} in total. The operation relied on a well-established and adaptable modus operandi, with the script adjusted according to the victim's profile. The case exemplifies the psychological sophistication of modern romance fraud: victims, often elderly or isolated, transferred their savings --- including pension payments --- in the belief they were supporting a loved one.

\subsection{Sextortion: The Jordan DeMay Case (March 2022)}

In March 2022, Jordan DeMay, a 17-year-old American high school football player from Michigan, was contacted via an Instagram account (``dani.robertts'') that appeared to belong to a young woman. The account had in fact been hacked and sold to a 22-year-old Nigerian man, who used it to coerce young men into sending explicit images of themselves, then threatened to distribute the images unless money was paid. DeMay died by suicide within hours of being targeted. The suspects, Samuel and Samson Ogoshi, were subsequently extradited to the United States and charged in federal court. The case became a landmark in debates about the criminal liability for sextortion-induced suicide and the cross-jurisdictional challenges of prosecuting online sexual coercion.

\begin{tcolorbox}[colback=yellow!5!white,colframe=orange!80!black]
The DeMay case exemplifies the link between IBA and assisted cyber suicide: sextortion does not merely constitute a financial or reputational harm, but in the most severe cases constitutes a precipitating cause of death. It also illustrates how the perpetration of sextortion does not require the offender to possess any intimate images in advance --- the scheme is engineered to create them.
\end{tcolorbox}

\subsection{Online Sexual Crimes Against Children: The Love Zone}

The Love Zone (TLZ) was one of the largest pedophile online communities in the world, with approximately \textbf{45,000 members}. The platform was used to disseminate pedophile practices and to sell, procure, and trade child sexual abuse material (CSAM). In 2014, Queensland Police identified the TLZ administrator, covertly assumed control of his account, and monitored the site for six months, exposing the forum's most active and prolific members in a classic law enforcement ``honeypot'' operation. The operation ultimately led to dozens of arrests across multiple jurisdictions; the most recent arrests linked to TLZ date to May and August 2023. One individual associated with the network was charged with 1,623 child abuse offences. The case highlights the tension between the investigative value of covertly operating illegal platforms (allowing broader identification of offenders) and the ethical and legal implications of temporarily permitting ongoing criminal activity.



% ========================================
\part{Cybercrime Typologies: Cyber-assisted Crimes}
% ========================================

\section{Definition}

Cyber-assisted crime is the third and least technology-dependent category in the Brenner--Wall trichotomy. These are traditional crimes in which the internet is used ``to assist with the organisation of a crime'' (Wall, 2005). The term \textit{cyber-assisted crime} was formally proposed by McGuire and Dowling (2013).

\begin{definition}[Cyber-assisted crime]
Traditional crimes in which the digital tool is \textbf{not critical}: the crime ``would still take place if the internet was removed'' (Wall, 2017). Technology is part of more complex offences or criminal activities, used to \textbf{assist pre-existing operations} or by moving part of the criminal business online (Wall, 2017; Di Nicola, 2022). The result is what Wall (2017) calls \textbf{cyber-assisted forms of organisation}.
\end{definition}

The key distinction from cyber-enabled crime is one of degree, not of kind: in cyber-assisted crime, digital tools intervene at specific phases of a criminal activity that is structurally independent of them. Remove the internet and the crime still happens; with the internet, it becomes more efficient, more scalable, and harder to disrupt.

\begin{remark}
The three categories --- cyber-dependent, cyber-enabled, cyber-assisted --- form a spectrum of digital integration, from crimes that cannot exist without technology, to crimes that merely use technology as an incidental organisational tool. Most real-world organised crime today operates somewhere along this spectrum rather than at a single fixed point.
\end{remark}


\section{Crime Scripts}

\subsection{The Concept}

To understand cyber-assisted crime, it is useful to adopt the analytical framework of \textbf{crime scripts} (Cornish, 1994). A crime script is a structured description of the sequential steps involved in committing a crime, identifying the actions taken, the resources required, and the decision points at which intervention is possible. Applied to cyber-assisted crime, crime scripts allow us to map exactly \textit{where} and \textit{how} digital tools intervene in an otherwise traditional criminal process.

Key features of the crime script approach:

\begin{itemize}
    \item Crime is treated as a process that can be ``unfolded into separate but related phases'' (Lavorgna, 2013).
    \item The influence of the digital society may be limited to \textbf{one or more phases} of the criminal activity, not necessarily the whole.
    \item The method was recommended by Cornish and Clarke (2002) specifically for investigating organised crime.
    \item The goal is to understand the role of the internet, the opportunities it creates, and the phases in which those opportunities occur --- providing law enforcement with \textbf{intervention points}.
\end{itemize}

\subsection{Example: Tag Writing (Cornish, 1994)}

The classic illustrative example of a crime script is graffiti tag writing. Each stage has a corresponding script action and a corresponding situational control measure that could prevent or disrupt it:

\begin{center}
\begin{tabular}{lll}
\hline
\textbf{Scene/Function} & \textbf{Script Action} & \textbf{Situational Control} \\
\hline
Preparation & Buy spray-can; find a good setting & Sales regulation; city paint-out program \\
Entry & Enter setting & Access control; entry/exit screening \\
Pre-condition & Loiter & Surveillance \\
Instrumental pre-condition & Select target & Remove target \\
Instrumental initiation & Approach target & Surveillance \\
Instrumental actualization & Reach target & Protective screens; legal target provided \\
Doing & Spray graffiti & Graffiti-resistant paint \\
Post-condition & Get away quietly & Moisture-activated alarm \\
Exit & Leave setting & Entry/exit screening \\
Doing (later) & ``Getting up'' & Rapid cleaning \\
\hline
\end{tabular}
\end{center}

The value of the script is not the description of any individual act but the systematic identification of \textit{all} the stages at which prevention could interrupt the sequence. The same logic applies to complex organised crimes.

\subsection{Application: Wildlife Trafficking (Lavorgna, 2013)}

Wildlife trafficking is a paradigmatic cyber-assisted crime: a multi-stage smuggling operation encompassing numerous activities (Wyatt, 2022), in which the internet plays a \textbf{crucial role specifically in the advertising and selling stages} --- making it easier to find potential buyers or sellers (Lavorgna, 2020) --- while physical logistics remain offline.

The crime script of wildlife trafficking (Lavorgna, 2013: 78) unfolds across seven stages:

\begin{enumerate}
    \item[\textbf{Stage 0}] Preparatory activities antecedent to the commission of wildlife trafficking.
    \item[\textbf{Stage 1}] Poaching, harvesting, or breeding of the animal or plant.
    \item[\textbf{Stage 2}] Intermediate passage through local middlemen and the domestic market.
    \item[\textbf{Stage 3}] Passage through regional middlemen and international traders.
    \item[\textbf{Stage 4}] Intermediate passage through local middlemen and the domestic market (at destination).
    \item[\textbf{Stage 5}] Distribution of the animal, plant, or product.
    \item[\textbf{Stage 6}] Activities directly consequential or subsequent to the trafficking activity.
\end{enumerate}

Lavorgna's detailed mapping of internet use across these stages (2013: 79--81) shows that digital tools appear across all stages as \textbf{Crime}, \textbf{Lifestyle}, and \textbf{Network} script categories. Selected examples:

\begin{itemize}
    \item \textbf{Preparation (Stage 0)}: Formation of criminal networks via email; use of dedicated forums to exchange information on potential buyers; use of Google Translate to post announcements in different languages; communication with other members of the network about relevant legislation.
    \item \textbf{Entry (Stage 0, 1, 4)}: Discussion with criminal peers to exchange information on potential new buyers; interaction with potential buyers via auction websites, general commercial websites, dedicated commercial websites, dedicated forums.
    \item \textbf{Pre-condition}: Maintenance of contacts with criminal network members (email, Skype); signalling via dedicated websites of recently discovered species that are rare and valuable; obtaining information on specific species (morphology, breeding); sharing information about the location of the animal/plant in its natural environment (dedicated blogs); agreement to encourage local middlemen who hesitate due to fear of prosecution; informing potential clients and middlemen of updated prices.
    \item \textbf{Instrumental pre-condition (Stages 3, 4)}: Agreement to exclude local middlemen who hesitate; search for new suppliers via emails and dedicated forums.
    \item \textbf{Instrumental actualization (Stages 3--5)}: Advertising via fairs and exhibitions (dedicated websites, forums); forwarding emails from suppliers who can fulfil orders; identifying the location of the animal via GPS and web mapping services (e.g.\ Google Maps); ordering/reserving specimens on dedicated commercial websites; arranging sales of the specimen to middlemen via international auction websites; discussing prices and meeting places in the physical world to close deals (email, domains); discussing via criminal network how to ship specimens in small packages with deceptive labels and then send them through regular mail; agreement to physically meet and exchange animals/products (private messaging on dedicated commercial websites, email); putting customers directly in contact with middlemen or suppliers; contacting customers to reassure them that, without documents, the price is significantly lower; discussing with potential customers the absence of required documentation; reassuring customers about the quality of the animal/plant/product on dedicated commercial websites and auction websites.
    \item \textbf{Doing (Stages 1, 4, and later)}: Auctioning/selling the animal/plant/product (dedicated commercial websites); selling on auction websites (e.g.\ eBay); agreement to send animals/plants/products via regular mail among middlemen; monitoring delivery via regular mail; payment via prepaid cards that can be recharged online.
    \item \textbf{Post-condition (Stages 4--6)}: Buying accessories and sharing videos and pictures in zoos or cases where wild animals are involved (dedicated forums, deep web); attempting to cover the criminal act by claiming that animals/plants/products have been imported in the past with CITES certificates but have been given to others (dedicated website).
\end{itemize}

This detailed mapping reveals that the internet penetrates every stage of the wildlife trafficking script, yet the crime itself --- the physical capture, transport, and delivery of animals --- remains independent of digital infrastructure. This is the hallmark of cyber-assisted crime.


\section{New Criminal Opportunities}

Modern technological tools have created \textbf{new criminal opportunities} for organised groups (Lavorgna, 2013, 2015a, 2015b). Understanding this dimension is essential for grasping the full impact of digital society on criminal activities, organised or otherwise.

Lavorgna identifies \textbf{five types of opportunities} available to criminal groups:

\begin{enumerate}
    \item \textbf{Opportunities for communication}: encrypted messaging platforms, dedicated forums, and anonymous communication tools allow criminal groups to coordinate at scale and across borders with greatly reduced risk of interception.

    \item \textbf{Management opportunities}: digital tools facilitate the logistics of criminal operations --- inventory management, financial transactions, coordination of supply chains --- making large and geographically dispersed operations manageable.

    \item \textbf{Organisational opportunities}: the internet enables new forms of criminal organisation, including loosely networked structures that are more resilient to law enforcement disruption than traditional hierarchical groups. Members may never meet physically.

    \item \textbf{Relational opportunities}: online platforms allow criminal groups to identify, vet, and recruit new members, build trust relationships with suppliers and buyers, and maintain ongoing relationships with co-offenders and clients across jurisdictions.

    \item \textbf{Promotional, marketing, persuasive, and fidelization opportunities}: criminal groups use the internet to advertise services and products to potential customers, build a reputation, manage reviews and feedback (as seen in dark web marketplaces), and retain clients through loyalty mechanisms analogous to legitimate e-commerce.
\end{enumerate}


\section{Cyber-assisted Criminal Activities}

Cyber-assisted crimes are divided into two broad clusters based on who commits them:

\subsection{Trafficking Activities of Collectivities}

These are multi-actor, multi-stage operations in which organised criminal groups use digital tools to assist pre-existing trafficking activities:

\begin{itemize}
    \item \textbf{Counterfeit medicine trafficking}: the internet facilitates the sourcing, advertising, and distribution of falsified pharmaceutical products, as explored in Project FAKECARE.
    \item \textbf{Smuggling of migrants and asylum seekers}: as documented in the Surf and Sound project, social media platforms and messaging apps are used at every stage of the smuggling process --- recruitment, route planning, real-time navigation, and feedback.
    \item \textbf{Trafficking for the purpose of exploitation} (particularly sexual): digital recruitment tools, fake job advertisements, and loverboy methods are used to identify and groom victims before physical exploitation occurs.
    \item \textbf{Terrorism and radicalisation}: the internet is used as a recruitment and propaganda tool, for coordination of attacks, and for dissemination of extremist ideology --- while the physical acts of violence remain independent of digital infrastructure.
    \item \textbf{Sale of counterfeit products or psychotropic substances}: online marketplaces (surface web, deep web, dark web) facilitate supply, demand, and payment for a wide range of illicit goods.
    \item \textbf{Illegal animal trafficking}: as illustrated by the wildlife crime script above.
    \item \textbf{Tobacco trade}: illicit tobacco trafficking uses online platforms for sourcing, ordering, and coordinating distribution networks.
\end{itemize}

\subsection{Criminal Activities of Individuals}

Beyond organised crime, individual criminals also exploit digital tools in an assistive capacity:

\begin{itemize}
    \item \textbf{Using the internet to learn how to commit or hide a crime}: online tutorials, forums, and dark web resources provide information on criminal techniques, from document forgery to evading forensic detection.
    \item \textbf{Using social media accounts to gather useful information about prospective victims}: OSINT (open-source intelligence) techniques applied to social media allow individual criminals to profile targets, identify vulnerabilities, and plan offences with a level of detail previously available only to professional investigators.
    \item \textbf{Crime-as-a-service}: the digital market allows individuals to outsource specific criminal tasks --- hacking, money mule recruitment, document forgery, phishing kit deployment --- without possessing the skills themselves, lowering the barrier to entry for complex criminal activity.
\end{itemize}


\section{Case Studies}

\subsection{Encrypted Criminal Communications: EncroChat (2020)}

EncroChat was an encrypted communication tool used on a large scale by organised criminal groups for criminal purposes. With an estimated \textbf{60,000 users}, the platform provided hardware-modified phones running a hardened operating system with features specifically designed to frustrate law enforcement: remote wipe capability, no GPS logging, and end-to-end encrypted messaging. In 2020, a joint investigation by French and Dutch authorities infiltrated the EncroChat network and was able to intercept, share, and analyse millions of communications in real time. The operation led to \textbf{over 6,500 arrests} across multiple countries and the seizure of close to \textbf{EUR 900 million}. Similar operations subsequently dismantled two further encrypted criminal communication networks: \textbf{Exclu} and \textbf{Sky ECC}.

The EncroChat case exemplifies the ``opportunities for communication'' that digital tools provide to organised crime. It also illustrates the law enforcement response: rather than blocking the communications, authorities chose to covertly monitor them, harvesting intelligence on criminal networks at an unprecedented scale. This raises significant questions about admissibility of evidence, bulk surveillance, and the rule of law that continue to be litigated in courts across Europe.

\subsection{Human Trafficking and Technology}

Technology plays an increasingly important role in trafficking in human beings (THB), particularly in the two key phases of the trafficking process: \textbf{recruitment} and \textbf{exploitation} of victims.

\paragraph{Recruitment phase.}
\begin{itemize}
    \item \textbf{Fake job advertisements}: attractive work conditions (high pay, foreign travel, hospitality work) are advertised on legitimate job platforms and social networks, luring victims into exploitative situations before the deception becomes apparent.
    \item \textbf{Loverboy method}: the trafficker cultivates a fake romantic relationship with the victim online, building emotional dependence before introducing conditions of exploitation. This method is particularly effective against young and socially isolated individuals.
\end{itemize}

\paragraph{Exploitation phase.}
Once recruited, digital tools are used to control victims (e.g.\ confiscating phones, monitoring communications), advertise victims' services online, collect payments, and launder proceeds.

The GRETA report (\textit{Online and technology-facilitated trafficking in human beings}) documents the systematic integration of digital tools into THB operations and the corresponding need for law enforcement responses adapted to the digital environment. A 2023 hackathon specifically targeted the problem of human traffickers luring victims online, bringing together technologists and law enforcement to develop detection tools.

\subsection{Illicit Online Markets}

Digital technologies are used to facilitate the sale of a wide range of illegal and contraband goods. Key categories include:

\begin{itemize}
    \item \textbf{Counterfeit products}: online platforms (including social media and dark web marketplaces) are used to advertise, sell, and ship counterfeit goods. An international operation by Europol shut down 12,526 websites offering counterfeit goods and pirated content, seizing 127,365 counterfeit products worth EUR 3.8 million.
    \item \textbf{Illegal dark web markets}: dedicated marketplaces on the dark web provide a full-service environment for buying and selling drugs, weapons, documents, and other contraband. Law enforcement on three continents arrested 288 dark web vendors in a major marketplace seizure, seizing EUR 50.8 million in cash and virtual currencies, 850 kg of drugs, and 117 firearms.
    \item \textbf{Wildlife and endangered species}: as analysed through the crime script framework above.
    \item \textbf{Cultural heritage artifacts}: the internet facilitates the illicit trade in looted or stolen cultural heritage objects. An international art trafficking operation led to 60 arrests and the recovery of over 11,000 objects.
    \item \textbf{Pharmaceuticals}: online markets, including end-to-end encrypted messaging platforms such as Telegram, are used to distribute both controlled and uncontrolled substances. A documented case involved an underground abortion pill market growing on Telegram as US states restricted access to abortion care, illustrating how legal and political context shapes the demand side of illicit online markets.
\end{itemize}

\begin{tcolorbox}[colback=yellow!5!white,colframe=orange!80!black]
Across all three case studies, a common pattern emerges: the criminal activity itself (trafficking, violence, physical delivery of goods) is not digital, but the internet provides the infrastructure for communication, recruitment, advertising, payment, and logistics. This is the operational core of cyber-assisted crime. Law enforcement responses must therefore target not only the physical crime but also the digital layers that enable it --- and doing so raises complex questions about privacy, jurisdiction, and the proportionality of surveillance.
\end{tcolorbox}



% ========================================
\part{AI, ML, and Measuring Cybercrime}
% ========================================

% --------------------------------------------------
\section*{Lecture A --- AI and ML and Cybercrime}
\addcontentsline{toc}{section}{Lecture A --- AI and ML and Cybercrime}
% --------------------------------------------------

\section{Human--Machine Interaction and Crime}

The literature on digital sociology emphasises the \textbf{relationship between the human being and the machine in the construction of the new digital society}, in a socio-technical vision. In a world where robots support or replace humans, and where artificial intelligence is gaining ground, a fundamental criminological question arises: are there, or will there be, increasingly \textbf{hybrid criminals} --- the result of close interaction between humans and computers?

King et al.\ (2020: 91) define phenomena of ``\textbf{crime perpetrated by artificial intelligence}'' in which the perpetrator is very difficult to circumscribe, and in which some questions arise spontaneously: \textit{Who commits the crime? A human agent? An artificial agent? Both of them?} These questions challenge the anthropocentric assumptions at the core of traditional criminology and criminal law.

\section{Artificial Intelligence: The Concept}

\begin{definition}[Artificial intelligence]
Artificial intelligence deals with \textbf{artificial agents capable of carrying out tasks that usually require intelligence and human intervention to be conducted}. These artificial agents are ``sufficiently informed'', ``intelligent'', ``autonomous and capable of carrying out morally relevant actions independently of the humans who created them'' (Floridi and Sanders, 2004: 351).
\end{definition}

The key criminological implication of this definition is the concept of \textbf{moral autonomy}: an AI system that can take actions independently of its creators creates a category of harmful conduct that sits outside the traditional offender--victim dyad. This forces criminology to reckon with questions of agency, responsibility, and attribution that existing frameworks were not designed to handle.

\section{AI and Criminology: Three Categories of Challenge}

AI introduces new challenges for crime and criminology along three axes:

\begin{enumerate}
    \item \textbf{Crime with artificial intelligence}
    \item \textbf{Crime against artificial intelligence}
    \item \textbf{Crime through artificial intelligence}
\end{enumerate}

\subsection{Crime with Artificial Intelligence}

AI can be a \textbf{powerful tool for criminal purposes}. As Brundage et al.\ (2018: 19--22) argue, AI in criminal hands is capable of:
\begin{itemize}
    \item \textbf{Expanding existing threats}: allowing criminals to scale attacks that were previously limited by human labour (e.g.\ automated phishing at scale).
    \item \textbf{Introducing entirely new threats}: creating categories of harm that were not previously possible.
    \item \textbf{Completely altering the characteristics of existing threats}: changing the speed, accuracy, cost, or detectability of a known crime.
\end{itemize}

The paradigmatic example is the \textbf{deepfake}: AI-generated audio-visual content that realistically impersonates a real person. Deepfakes can be used for disinformation (fabricating statements by public figures), financial fraud (voice cloning to impersonate executives), and image-based sexual abuse (generating non-consensual intimate imagery using a real person's likeness). As noted in Part IV with the Crosetto case (2025), deepfake technology has already been operationally deployed in large-scale fraud.

\subsection{Crime against Artificial Intelligence}

Here we are dealing with \textbf{attacks that exploit the vulnerabilities of an AI system to manipulate or ``hypnotise'' the system itself}. The AI system is the victim of the attack, but the ultimate objective is typically to weaponise the system against third parties.

Key attack vectors include:

\begin{itemize}
    \item \textbf{Training data poisoning}: the attacker corrupts the data used to train an AI model, causing the model to learn systematically distorted patterns. If the poisoned model is then deployed in a criminal justice, medical, or financial context, its decisions will be systematically biased in ways that benefit the attacker.
    \item \textbf{Adversarial inputs}: carefully crafted inputs that cause an AI system to misclassify or misbehave at inference time --- for example, fooling an image classifier or a facial recognition system.
    \item \textbf{Model reprogramming}: manipulating an AI system to commit criminal acts, behaving as an autonomous criminal agent.
\end{itemize}

The paradigmatic example is \textbf{Microsoft's Tay chatbot} (2016): a conversational AI deployed on Twitter that was manipulated within 24 hours through coordinated user input, causing it to generate racist and offensive content at scale. This is a case where coordinated social engineering constituted a form of ``crime against AI'' that turned the platform's own technology into a vector of harm --- with no clear individual perpetrator.

\subsection{Crime through Artificial Intelligence}

In this category AI becomes an \textbf{intermediary agent}: a non-human actor that stands between the human criminal and the criminal act, creating a layer of separation that complicates attribution and detection.

The clearest example is a \textbf{bot purchasing drugs on the dark web} instead of a natural person. The human criminal programmes the bot with intent and provides resources, but the actual criminal transaction is carried out autonomously by the AI. In this model, artificial intelligence acts as a \textbf{criminal shield} --- a screen between the offender and the offence that exploits the inability of legal systems to assign criminal liability to non-human agents.

\section{AI and Crime Areas}

AI-facilitated crime is not evenly distributed across crime types. The principal affected areas are:

\begin{itemize}
    \item \textbf{Identity theft and fraud, impersonation, and online fraud generally}: AI enables the creation of synthetic identities, voice and face cloning for impersonation, and automated phishing campaigns personalised at scale.
    \item \textbf{Financial crime}: algorithmic manipulation of markets, AI-driven money laundering through complex transaction patterns, and automated detection evasion.
    \item \textbf{Drug trafficking and trading}: AI bots as purchasing agents, algorithmic pricing and inventory management on dark web markets.
    \item \textbf{Harassment}: AI-generated hate speech, automated pile-ons, and the use of AI to identify and profile harassment targets.
    \item \textbf{Sex crimes}: deepfake pornography, AI-assisted grooming (chatbots maintaining grooming conversations at scale), and automated generation and distribution of child sexual abuse material.
\end{itemize}

\section{AI-enabled Future Crimes: Risk Matrix}

Caldwell, Andrews, Tanay et al.\ (\textit{Crime Science}, 9:14, 2020) systematically assessed AI-enabled future crimes along two dimensions:
\begin{itemize}
    \item \textbf{Harm or profit potential} (low to high)
    \item \textbf{Difficulty of defeating} (easier to harder)
\end{itemize}

The resulting risk matrix classifies AI-enabled crimes as follows:

\paragraph{High risk (high harm, hard to defeat):}
\begin{itemize}
    \item Audio/video impersonation (deepfakes)
    \item Tailored phishing (AI-personalised social engineering)
    \item AI-authored fake news
    \item Disrupting AI-controlled systems
    \item Large-scale blackmail
    \item Market bombing
\end{itemize}

\paragraph{Medium risk:}
\begin{itemize}
    \item Learning-based cyber attacks
    \item Tricking face recognition systems
    \item Online eviction
    \item Data poisoning
\end{itemize}

\paragraph{Low risk (low harm or easier to defeat):}
\begin{itemize}
    \item Bias exploitation
    \item AI-authored fake reviews
    \item Burglar bots
    \item AI-assisted stalking
    \item Forgery
    \item Evading AI detection
\end{itemize}

The matrix is useful for prioritising defensive investment: high-risk threats that are also difficult to defeat require urgent attention, while lower-risk or more easily defeated threats can be addressed with proportionally lower resources.

\section{Cyborg-Crime}

Working on botnets --- networks of infected computers controlled by a single operator, increasingly used to commit cybercrime --- van der Wagen and Pieters (2015) coined the concept of \textbf{cyborg-crime}: online crime perpetrated by a \textbf{cyborg actor}, and proposed the use of \textbf{Actor-Network Theory} (ANT) as a framework for interpreting forms of delinquency committed by actors who are \textit{neither completely human, nor completely driven by a machine}.

Their core argument: it is not possible to understand the nature of crimes carried out through botnets if we continue to use an \textbf{anthropocentric criminology} that sees the human being as the main force behind criminal behaviour. ``Rather than considering a botnet as a technological tool, an individual (opportunistic) crime or a resource on the criminal market, we treat a botnet as a \textbf{hybrid crime network}, a crime that derives from mutual cooperation and human/technology interaction'' (ibid.: 579).

The cyborg-crime concept bridges the previous three categories (with/against/through AI) and points towards a broader theoretical reconceptualisation of criminal agency in the digital age. It challenges criminology to develop frameworks capable of assigning responsibility, assessing harm, and designing interventions in contexts where the boundary between human and machine agency is systematically blurred.


% --------------------------------------------------
\section*{Lecture B --- Measuring Cybercrime: Data Sources and Challenges}
\addcontentsline{toc}{section}{Lecture B --- Measuring Cybercrime: Data Sources and Challenges}
% --------------------------------------------------

\section{Overview: Crime Data Sources}

Traditional crimes can be measured through three main data sources:

\begin{enumerate}
    \item \textbf{Official crime statistics}: quantitative data produced by the criminal justice system.
    \item \textbf{Victimisation surveys}: data collected directly from the population about their experience of crime.
    \item \textbf{Self-report surveys}: data collected from individuals about their own criminal or deviant behaviour.
\end{enumerate}

Cybercrime, however, is extremely difficult to measure and to reconcile across these sources. Each instrument faces specific challenges when applied to digital crime, compounding the already severe \textbf{dark figure of crime} problem.

\section{Official Crime Statistics}

\subsection{Definition and History}

\begin{definition}[Official crime statistics]
Quantitative data on crime recorded by state agencies, primarily police, prosecutors, and courts. Originated in France in the early 19th century, initially focused on judicial and penitentiary statistics and later extended to police-recorded data.
\end{definition}

Main administrative sources at the supranational level:
\begin{itemize}
    \item \textbf{Interpol}: collected comparative data until 2006.
    \item \textbf{UN Crime Trends Survey (CTS)}: collects standardised crime statistics from member states.
    \item \textbf{European Sourcebook of Crime and Criminal Justice Statistics}: compiles cross-national data from European countries.
    \item \textbf{Eurostat database}: produces harmonised EU-level crime statistics.
\end{itemize}

Key structural problems:
\begin{itemize}
    \item The \textbf{dark figure of crime}: a large proportion of crimes are never reported to police, never recorded, or lost at subsequent stages of the criminal justice process. Official statistics therefore systematically undercount actual crime.
    \item \textbf{No detailed information on the criminal event, the victim, or the offender}: administrative records capture only what is legally and procedurally relevant, stripping out the contextual information needed for criminological analysis.
\end{itemize}

\subsection{Official Statistics and Cybercrime}

Official statistics on cybercrime are \textbf{extremely limited}, for several structural reasons:

\begin{itemize}
    \item Cybercrimes --- especially cyber-enabled ones --- are in most cases \textbf{registered as traditional crimes}. A phishing fraud is recorded as fraud; a sextortion case may be recorded as blackmail. The digital dimension is lost.
    \item There is a \textbf{lack of uniform registration} of cybercrime in the police registration systems of most Western countries: no agreed taxonomy, no mandatory cybercrime flag, no standardised coding.
    \item Information concerning the digital nature of a crime may be present in \textbf{unstructured qualitative text fields} in police or court registration systems --- but extracting it is a labour-intensive task not amenable to standard statistical processing.
\end{itemize}

\paragraph{Notable exceptions:}
\begin{itemize}
    \item \textbf{FBI Internet Crime Complaint Center (IC3)}: a dedicated reporting portal for internet crime in the United States, producing annual reports on cybercrime complaints and losses.
    \item \textbf{Uniform Crime Report Survey of Canada}: includes specific cybercrime categories.
    \item \textbf{European Sourcebook of Crime and Criminal Justice Statistics}: publishes cross-national records of cyber-enabled fraud and, for some countries, cyber-dependent crime.
\end{itemize}

\paragraph{Methodological response --- machine learning:}
To overcome the problem of cybercrime information buried in unstructured text fields, two research approaches have been developed:
\begin{itemize}
    \item Sessink (University of Twente): ``\textit{Using Machine Learning to Detect ICT in Criminal Court Cases}'' --- applying NLP classifiers to court documents to identify cases with an ICT dimension.
    \item Van der Laan and Tollenaar: ``\textit{Text Mining for Cybercrime in Registrations of the Dutch Police}'' --- using text mining to flag cybercrime-relevant entries in police records that would otherwise be invisible to quantitative analysis.
\end{itemize}

\section{Victimisation Surveys}

\subsection{Definition}

\begin{definition}[Victimisation survey]
A technique that involves asking a representative sample of the population: (1) whether and of what crimes they have been victims in a given period of time prior to the interview; and (2) if so, further details about the incident (where it occurred, at what time, and how).
\end{definition}

The key advantage over official statistics is that victimisation surveys capture crimes that were never reported to police, directly addressing the dark figure of crime.

\subsection{The US National Crime Victimization Survey (NCVS)}

The NCVS is the foundational instrument of modern victimisation research. Key milestones:

\begin{itemize}
    \item \textbf{First pilot studies in the 1960s}: conducted on 10,000 people, respondents reported \textbf{double the crime rate recorded in official statistics} --- a landmark finding that demonstrated the scale of the dark figure and justified the survey approach.
    \item \textbf{Since 1972}: the NCVS has been continuously investigating the frequency, characteristics, and consequences of victimisation in the United States.
    \item \textbf{Coverage}: crimes against the person (rape and sexual assault, robbery, assault, bag snatching/pickpocketing) and against property (theft, burglary, motor vehicle theft).
\end{itemize}

The NCVS allows researchers to:
\begin{enumerate}
    \item \textbf{Estimate the total number of crimes} committed against the population as a whole --- consistently showing a far higher figure than official statistics (the gap between the NCVS and the Uniform Crime Report (UCR) is visually striking, with NCVS figures for violent crime running at several times the UCR count for the same years).
    \item \textbf{Investigate differential victimisation}: analyse the levels of victimisation to which different segments of the population are exposed, based on victim characteristics (age, gender, ethnicity, marital status, income, education). Example: males, Black individuals, and young people up to age 24 are more at risk of violent crime than women, white individuals, and those over 24.
    \item \textbf{Provide details about the perpetrator and the events}: capturing information on offender demographics, victim--offender relationship, and circumstances of the offence that official statistics do not record.
\end{enumerate}

\paragraph{Impact of the NCVS:}
The success of the NCVS prompted the development of similar surveys at both international and national levels:
\begin{itemize}
    \item \textbf{International}: the \textbf{International Crime Victims Survey (ICVS)}.
    \item \textbf{National}: Netherlands (Victims of Crime, 1974), United Kingdom (British Crime Survey, 1982), Italy (1997).
\end{itemize}

\subsection{Victimisation Surveys and Cybercrime}

The application of victimisation surveys to cybercrime is still developing. A systematic review by Reep-van den Bergh and Junger (\textit{Crime Science}, 7:5, 2018) --- ``\textit{Victims of cybercrime in Europe: a review of victim surveys}'' --- provides the most comprehensive picture to date:

\begin{itemize}
    \item \textbf{Objective}: to map cyber-victimisation patterns in Europe and provide estimates of annual prevalence rates (number of victims per population) for six cybercrimes.
    \item \textbf{Methodology}: definition of survey inclusion criteria; database and internet searches; grey literature; contacts with national statistical offices.
    \item \textbf{Results}: nine surveys identified across European countries; highly heterogeneous annual prevalence rates depending on the crime type:
    \begin{itemize}
        \item 1--3\% for online shopping scams.
        \item 3\% of the population is a victim of cyberbullying.
    \end{itemize}
\end{itemize}

The wide variation in estimates across surveys --- arising from differences in operationalisation, reference periods, sampling frames, and question wording --- underlines how immature the field of cybercrime measurement still is. Comparability across surveys remains a major methodological challenge.

\section{Self-Report Surveys}

\subsection{Definition and Purpose}

Self-report surveys are instruments in which respondents are asked about instances in which they have been \textbf{actively involved in crime or deviant behaviour} --- i.e., as offenders rather than victims. They are particularly important for:

\begin{itemize}
    \item Measuring \textbf{juvenile delinquency}, where official statistics are especially unreliable due to selective enforcement and diversion.
    \item Developing and testing \textbf{theories of criminal and deviant behaviour}, since they provide data on the prevalence and frequency of offending in the general population rather than only among those caught by the system.
\end{itemize}

The key international self-report instrument is the \textbf{International Self-Report Delinquency Study (ISRD)}.

\subsection{Self-Report Surveys and Cybercrime}

A growing body of self-report surveys incorporates cybercrime measures. Three key instruments (Buil-Gil, Trajtenberg and Aebi, 2025) are:

\begin{center}
\begin{tabular}{llllll}
\hline
\textbf{Name} & \textbf{Acronym} & \textbf{Country} & \textbf{Frequency} & \textbf{Round} & \textbf{Sample} \\
\hline
Offending, Crime & OCJS & UK & Annual, & 2006 & 5,354 respondents \\
and Justice Survey & & & 2003--2006 & & aged 10--29 \\
Youth Delinquency & YDS & Netherlands & Quinquennial & 2015 & 1,471 respondents \\
Survey & & & & & aged 12--23 \\
International & ISRD4 & Cross- & 1991/91, & 2021/22 & Varies by country \\
Self-Report & & national & 2006/08, & & (aged 13--17) \\
Delinquency Study & & & 2012/19, & & \\
 & & & 2021/22 & & \\
\hline
\end{tabular}
\end{center}

\paragraph{Measures of cybercrime offending used across these surveys (Table 6, Buil-Gil et al., 2025):}

\begin{itemize}
    \item \textbf{Malware}: OCJS --- ``used internet to send viruses on purpose to other computers''; YDS --- ``intentionally sent out viruses through e-mail or over the internet''.
    \item \textbf{Hacking}: OCJS --- ``used internet to hack into other computers? by hacking we mean using a computer to illegally get access to another computer's files''; YDS --- ``logged on to somebody else's computer, email or social media account without informed consent''; ISRD4 --- ``hacked or broken into a private account or computer to acquire data, get control of an account, or destroy data''.
    \item \textbf{DoS}: YDS --- ``tried to disrupt a website or email account by sending out large amounts of data''.
    \item \textbf{Digital piracy}: OCJS --- ``used internet to download software, music or films that you knew to be pirated or unauthorised''.
    \item \textbf{Online banking/shopping fraud}: OCJS --- ``bought anything over the internet using payment card or card details that did not belong to you, without the card owner's permission''; YDS --- ``sold something through internet, but not sending out goods after receiving payment''; ISRD4 --- ``used internet, e-mail or social media to dupe others (like phishing, selling worthless or illegal things, etc.) to make money''.
    \item \textbf{Hate speech}: ISRD4 --- ``sent hurtful messages or comments on social media about someone's race, ethnicity or nationality, religion, gender identity, etc.''.
    \item \textbf{Online harassment}: OCJS --- ``sent email message to someone in order to harass, scare or threaten them''; YDS --- ``threatened someone through text messages, e-mails or in chat boxes''.
\end{itemize}

\begin{tcolorbox}[colback=yellow!5!white,colframe=orange!80!black]
A recurring observation across all three measurement instruments: the dark figure of cybercrime is almost certainly larger than for traditional crime. Victims do not recognise they have been victimised; reporting rates are extremely low; official categories do not map onto cybercrime; and offenders who commit cybercrime often do not perceive their behaviour as criminal. Measuring the true extent of cybercrime remains one of the most significant methodological challenges in contemporary criminology.
\end{tcolorbox}


\end{document}
